% sections/07_augmentation.tex  –  Data Augmentation (Stage 3)
% \input{} into main.tex; do NOT wrap in \documentclass or \begin{document}

\section{Data Augmentation}
\label{sec:augment}

Augmentation is \textbf{Stage~3} of the data-conditioning pipeline
introduced in \Cref{sec:preproc}.  Its purpose is to synthetically
expand the training set so that models — particularly deep networks
with high parameter counts (\Cref{sec:overfit}) — see greater input
variety during optimisation.  Two principles govern every augmentation
applied in this project:

\begin{enumerate}
  \item \textbf{Physics-awareness.}  Every transformation must have a
        credible physical or instrumental justification.
        $S_{11}$ traces encode the electromagnetic response of a sensor
        in a specific physical state; arbitrary perturbations that
        violate Maxwell's equations or produce non-physical spectral
        shapes do not improve real-world robustness.
  \item \textbf{Training-time only.}  Augmented samples are generated
        on-the-fly (or pre-generated and cached) for the training
        partition exclusively.  Validation and test data are never
        augmented; doing so would corrupt the evaluation protocol.
\end{enumerate}

\begin{tcolorbox}[
  colback=orange!5!white,
  colframe=orange!75!black,
  title={\textbf{Warning — Plausibility vs.\ Diversity}},
  fonttitle=\bfseries
]
Augmentations that create visually diverse but physically implausible
samples can inflate the apparent dataset size without improving
real-world generalisation.  Each augmentation must have a credible
physical or instrumental justification.  When in doubt, ablate: train
with and without the augmentation and compare validation performance.
\end{tcolorbox}

%% ------------------------------------------------------------------
\subsection{Additive Gaussian Noise}
\label{sec:augment:noise}

\paragraph{Physical justification.}
Every VNA measurement contains additive thermal and receiver noise.
The noise floor depends on the instrument's dynamic range, IF bandwidth,
and averaging factor.  Injecting Gaussian noise simulates repeated
measurements of the same physical state under slightly different
instrument conditions.

\paragraph{Implementation.}
For each training sample $\mathbf{x} \in \mathbb{R}^{4001}$ (magnitude
channel), generate a perturbed copy:
\begin{equation}
  \tilde{\mathbf{x}} = \mathbf{x} + \boldsymbol{\varepsilon},
  \qquad
  \varepsilon_i \sim \mathcal{N}(0,\,\sigma_{\text{noise}}^{2}),
  \label{eq:gauss_noise}
\end{equation}
where $\sigma_{\text{noise}}$ is chosen to match a realistic
signal-to-noise ratio (SNR).  For a modern benchtop VNA with a dynamic
range of \SI{80}{\decibel} or more, noise levels on the order of
\SIrange{0.01}{0.1}{\decibel} per point are representative.  The same
procedure is applied independently to the phase channel with an
appropriate phase-noise variance (typically a few hundredths of a
degree).

\paragraph{Risk level: \textbf{Low}.}
Additive instrument noise is a well-understood corruption mechanism.
The augmentation is compatible with both Route~A and Route~B.

%% ------------------------------------------------------------------
\subsection{Frequency-Local Smooth Perturbation}
\label{sec:augment:jitter}

\paragraph{Physical justification.}
Minor calibration drift between measurement sessions can introduce
smooth, low-frequency distortions in the magnitude and phase responses.
These manifest as gentle undulations superimposed on the true spectral
shape rather than as point-wise noise.

\paragraph{Implementation.}
Generate a smooth perturbation curve by sampling a small number
($K = 5$--$10$) of random control points along the frequency axis,
drawing their amplitudes from $\mathcal{N}(0,\,\sigma_{\text{jitter}})$,
and interpolating with a cubic spline to produce a continuous
perturbation vector $\boldsymbol{\delta} \in \mathbb{R}^{4001}$.
The perturbed trace is $\tilde{\mathbf{x}} = \mathbf{x} +
\boldsymbol{\delta}$.

The amplitude $\sigma_{\text{jitter}}$ must be small enough that the
resulting trace remains within the class-aggregated envelope identified
during outlier screening (\Cref{sec:preproc}).

\paragraph{Risk level: \textbf{Low--Medium}.}
The perturbation is physically motivated but the amplitude must be
bounded carefully; excessive jitter can blur inter-class boundaries.

%% ------------------------------------------------------------------
\subsection{Magnitude Scaling}
\label{sec:augment:magscale}

\paragraph{Physical justification.}
Variations in cable loss, connector torque, and adapter insertion loss
introduce a quasi-constant multiplicative offset (in log-magnitude, an
additive offset in dB) across the frequency band.  Simulating this
variation makes the classifier robust to test-time setup differences.

\paragraph{Implementation.}
Apply a per-sample random gain $g$ drawn uniformly from
$[\SI{-1.5}{\decibel},\,\SI{+1.5}{\decibel}]$:
\begin{equation}
  \tilde{x}_i = x_i + g, \qquad g \sim \mathcal{U}(-1.5,\;1.5),
  \label{eq:mag_scale}
\end{equation}
where $x_i$ is the magnitude at frequency bin~$i$ in dB.
The $\pm\SI{1.5}{\decibel}$ range is representative of typical cable
and connector variability in X-band (\SIrange{8}{12}{\giga\hertz})
setups.

\paragraph{Risk level: \textbf{Low}.}
A uniform gain offset preserves the spectral \emph{shape}, which
carries the class-discriminative information.  Route~A and Route~B
compatible.

%% ------------------------------------------------------------------
\subsection{Phase Perturbation}
\label{sec:augment:phase}

\paragraph{Physical justification.}
Electrical length changes due to cable flexure or temperature drift
introduce a nearly linear phase offset
$\Delta\phi(f) \approx \alpha + \beta f$ across the band.
Small additional smooth fluctuations simulate connector phase
instability.

\paragraph{Implementation.}
For each training sample, generate a random affine phase offset:
\begin{equation}
  \tilde{\phi}_i = \phi_i + \alpha + \beta\, f_i,
  \label{eq:phase_perturb}
\end{equation}
where $\alpha \sim \mathcal{U}(-5^{\circ},\,5^{\circ})$ is a constant
offset and $\beta$ is drawn such that the total additional phase slope
across the \SI{4}{\giga\hertz} band does not exceed
$\pm\SI{10}{\degree}$.  Optionally, a smooth spline perturbation
(as in \Cref{sec:augment:jitter}) may be added to simulate non-linear
phase instability, but the total perturbation must remain small enough
to avoid creating non-physical discontinuities or phase wrapping
artefacts.

\paragraph{Risk level: \textbf{Medium}.}
Phase perturbations are physically justified but require careful
bounding.  Excessive phase noise can create unrealistic spectral
features, particularly near resonances where phase transitions are
steep.

%% ------------------------------------------------------------------
\subsection{Frequency-Axis Shift}
\label{sec:augment:freqshift}

\paragraph{Physical justification.}
A rigid shift of the frequency axis simulates a change in the VNA's
start/stop frequencies or a systematic frequency-reference error.  This
augmentation is \textbf{only justified} if the sensor acquisition
protocol includes variable centre-frequency settings or IF bandwidth
changes that alter the effective frequency grid.

\paragraph{Implementation (conditional).}
Shift the entire trace by $\Delta f$ frequency bins (with edge padding):
\begin{equation}
  \tilde{x}_i = x_{i + \Delta f},
  \qquad \Delta f \sim \mathcal{U}(-N_{\text{shift}},\;N_{\text{shift}}),
  \label{eq:freq_shift}
\end{equation}
where $N_{\text{shift}}$ is a small integer (e.g., $1$--$5$ bins,
corresponding to \SIrange{1}{5}{\mega\hertz}).

\paragraph{Risk level: \textbf{Medium--High}.}
If the acquisition protocol uses a fixed frequency plan (as appears to
be the case in this dataset), frequency-axis shifting has no physical
basis and should \textbf{not} be applied.  It is listed here for
completeness and should be reconsidered only if the data provider
confirms variable acquisition settings.

%% ------------------------------------------------------------------
\subsection{Mixup Augmentation}
\label{sec:augment:mixup}

Mixup~\cite{zhang2018mixup} creates new training samples as convex
combinations of existing ones:
\begin{equation}
  \tilde{\mathbf{x}} = \lambda\,\mathbf{x}_a + (1-\lambda)\,\mathbf{x}_b,
  \qquad
  \tilde{y} = \lambda\,y_a + (1-\lambda)\,y_b,
  \qquad
  \lambda \sim \text{Beta}(\alpha,\alpha),
  \label{eq:mixup}
\end{equation}
where $(\mathbf{x}_a, y_a)$ and $(\mathbf{x}_b, y_b)$ are two
randomly selected training samples and $\alpha \in (0,1]$ controls the
interpolation strength.

\paragraph{Risk assessment: \textbf{High}.}
Mixup was designed for natural-image classification where interpolated
images remain perceptually plausible.  For $S_{11}$ traces, an
interpolation between a \texttt{dry} trace and an \texttt{ice} trace
produces a spectrum that does not correspond to any physical sensor
state.  Furthermore, the soft label $\tilde{y}$ is uninterpretable in
the sensor context: there is no meaningful ``$40\,\%$~dry,
$60\,\%$~ice'' state.

\paragraph{Recommendation.}
Mixup should be used \emph{only} for intra-class interpolation (both
$\mathbf{x}_a$ and $\mathbf{x}_b$ from the same class) or avoided
entirely.  If employed inter-class, the benefit must be demonstrated
via careful ablation against a non-Mixup baseline, and the
interpretability implications must be documented.

%% ------------------------------------------------------------------
\subsection{SMOTE Prohibition}
\label{sec:augment:smote_prohibition}

\begin{tcolorbox}[
  colback=red!5!white,
  colframe=red!75!black,
  title={\textbf{SMOTE Is Prohibited for This Dataset}},
  fonttitle=\bfseries
]
SMOTE (Synthetic Minority Over-sampling Technique) and its variants
(ADASYN, Borderline-SMOTE) are \textbf{not permitted} in this pipeline.
The physical justification is threefold:
\begin{enumerate}
  \item \textbf{Non-physical interpolation.}  SMOTE creates synthetic
        samples by linear interpolation between existing minority-class
        neighbours.  An interpolation between a \texttt{dry} trace and
        an \texttt{ice} trace produces a spectrum that corresponds to no
        real dielectric state.  The three sensor surface states
        (dry, wet, ice) are \emph{discrete physical conditions}, not
        points on a continuous manifold.
  \item \textbf{Unnecessary given mild imbalance.}  The class ratio
        ($1.20:1:1$) is well within the range addressable by
        class-weighted loss functions and stratified cross-validation,
        without resorting to synthetic sample generation.
  \item \textbf{Risk of training on artefacts.}  Synthetic spectra that
        violate Maxwell's equations may teach the model decision
        boundaries that generalise poorly to real-world measurements.
\end{enumerate}
\end{tcolorbox}

This prohibition is consistent with the Mixup restriction above and
applies to all pipeline stages.  For class-balance adjustment, use
\texttt{class\_weight='balanced'} in the loss function or stratified
sampling strategies instead.

%% ------------------------------------------------------------------
\subsection{Image-Domain Augmentations for Route~B}
\label{sec:augment:image}

Route~B encodes $S_{11}$ traces as 2D images (e.g., magnitude heatmaps,
spectrograms, or recurrence plots).  Standard computer-vision
augmentations become available, but their applicability must be evaluated
against the underlying physics:

\begin{itemize}
  \item \textbf{Horizontal flip} (frequency-axis reversal):
        reverses the order of frequency bins.  This has no physical
        interpretation — $S_{11}$ at \SI{8}{\giga\hertz} and
        $S_{11}$ at \SI{12}{\giga\hertz} are not interchangeable —
        but may act as a form of regularisation if the classifier does
        not rely on absolute frequency position.
        \emph{Risk: Medium.}

  \item \textbf{Vertical flip} (magnitude-axis reversal):
        inverts the magnitude scale, creating physically impossible
        positive-dB returns.  \emph{Risk: High.  Not recommended.}

  \item \textbf{Random rotation}:
        rotates the image by an arbitrary angle.  A magnitude-vs.\
        -frequency plot rotated by even a few degrees mixes frequency
        and magnitude axes, destroying the physical meaning of both.
        \emph{Risk: Very high.  Not recommended.}

  \item \textbf{Random crop / resize}:
        crops a sub-region and resizes to the original dimensions.
        Acceptable only if the crop preserves the full frequency range;
        cropping out a resonant feature removes class-discriminative
        information.
        \emph{Risk: Medium.}

  \item \textbf{Colour jitter / brightness perturbation}:
        for single-channel greyscale images, brightness perturbation is
        analogous to magnitude scaling (\Cref{sec:augment:magscale})
        and shares the same justification and risk level.
        \emph{Risk: Low.}
\end{itemize}

%% ------------------------------------------------------------------
\subsection{Augmentation Summary Table}
\label{sec:augment:table}

\Cref{tab:augment_summary} consolidates all augmentation techniques,
their physical justification, risk level, and route compatibility.

\begin{table}[ht]
\centering
\caption{Summary of augmentation techniques for $S_{11}$ trace data.}
\label{tab:augment_summary}
\small
\renewcommand{\arraystretch}{1.15}
\begin{tabular}{p{3.2cm} p{3.6cm} p{1.4cm} p{1.7cm}}
\toprule
\textbf{Technique} &
\textbf{Physical justification} &
\textbf{Risk} &
\textbf{Route} \\
\midrule
Additive Gaussian noise &
  VNA thermal / receiver noise floor &
  Low &
  A, B \\
Smooth freq.-local jitter &
  Calibration drift between sessions &
  Low--Med &
  A, B \\
Magnitude scaling ($\pm$\,dB) &
  Cable loss / connector variation &
  Low &
  A, B \\
Phase perturbation &
  Electrical-length / temperature drift &
  Medium &
  A \\
Frequency-axis shift &
  Variable centre-freq.\ (conditional) &
  Med--High &
  A \\
Mixup (inter-class) &
  None (mathematical regulariser) &
  High &
  A, B \\
Mixup (intra-class) &
  Interpolation within physical state &
  Medium &
  A, B \\
Horizontal flip (image) &
  None (regulariser) &
  Medium &
  B \\
Vertical flip (image) &
  None (non-physical) &
  High &
  B \\
Random rotation (image) &
  None (non-physical) &
  Very high &
  B \\
Random crop / resize &
  Sub-band analysis (if full range kept) &
  Medium &
  B \\
Brightness perturbation &
  Analogous to magnitude scaling &
  Low &
  B \\
\bottomrule
\end{tabular}
\end{table}

%% ------------------------------------------------------------------
\subsection{Augmentation Budget and Ablation}
\label{sec:augment:ablation}

The number of augmented copies per original sample (the
\emph{augmentation factor}) must be chosen with care.  A factor of
$k$ transforms the effective training-set size from $N_{\text{train}}$
to $(1+k)\,N_{\text{train}}$, but the information content does not
scale linearly because augmented samples are correlated with their
parents.  Diminishing returns set in rapidly; empirical guidance from
the small-sample regime suggests $k \in [2,\,5]$ as a reasonable
starting range.

\paragraph{Ablation protocol.}
For every augmentation technique that is included in the final pipeline,
a controlled ablation must be performed:
\begin{enumerate}
  \item Train the model \emph{without} the augmentation and record
        validation performance (mean $\pm$ std over repeated CV).
  \item Train with the augmentation under identical hyper-parameters and
        compare.
  \item Retain the augmentation only if it produces a statistically
        meaningful improvement (e.g., paired $t$-test or Wilcoxon
        signed-rank test on fold-level metrics).
\end{enumerate}

This evidence-based approach prevents the accumulation of augmentations
that individually seem harmless but collectively distort the training
distribution away from the true data manifold.
